\documentclass[a4paper, UTF-8]{article}
\usepackage[margin=1in]{geometry}
\usepackage{algorithm}
\usepackage{algpseudocode}
\usepackage{amsmath}
\usepackage{amsfonts}
\usepackage{amsthm}

\usepackage[utf8]{inputenc} % allow utf-8 input
\usepackage[T1]{fontenc}    % use 8-bit T1 fonts
\usepackage{hyperref}       % hyperlinks
\usepackage{url}            % simple URL typesetting
\usepackage{booktabs}       % professional-quality tables
\usepackage{amsfonts}       % blackboard math symbols
\usepackage{nicefrac}       % compact symbols for 1/2, etc.
\usepackage{microtype}      % microtypography
\usepackage{xcolor}         % colors

\theoremstyle{plain}
\newtheorem{theorem}{\indent Theorem}[section]
\newtheorem{lemma}[theorem]{\indent Lemma}
\newtheorem{assumption}{\indent Assumption}[section]
\newtheorem{note}[theorem]{\indent Notation}
\newtheorem{proposition}[theorem]{\indent Proposition}
\newtheorem{corollary}[theorem]{\indent Corollary}
\newtheorem{definition}{\indent Definition}[section]
\newtheorem{example}{\indent Example}[section]
\newtheorem{remark}{\indent Remark}[section]
\newenvironment{solution}{\begin{proof}[\indent\bf Solution]}{\end{proof}}
\renewcommand{\proofname}{\indent\bf Proof}

\begin{document}

\title{\bf{Assignment 2}}

\author{\bf{Yanqiao Chen}\\ 
  Department of Computer Science and Engineering\\
  Southern University of Science and Technology\\
  \texttt{12412115@mail.sustech.edu.cn}
  }


\maketitle

\section{Page 27, Problem 7}

\begin{proof}
From the formula $P(A\cup B) = P(A) + P(B) - P(AB)$ and $P(A \cup B) \leq 1$, we have:
$$
P(A) + P(B) - P(AB) \leq 1
$$
Rearranging the inequality, we get:
$$
P(AB) \geq P(A) + P(B) - 1
$$
This complete the whole proof. $\square$
\end{proof}

\section{Page 27, Problem 8}

\begin{proof}

We prove this by mathematical induction. 

When $n=2$, we have:
$$
P(A_1 \cup A_2) = P(A_1) + P(A_2) - P(A_1 A_2) \leq P(A_1) + P(A_2)
$$

Assume that the inequality holds for $n=k$, that is:
$$
P\left(\bigcup_{i=1}^k A_i\right) \leq \sum_{i=1}^k P(A_i)
$$

Then for $n=k+1$, we have:

\begin{align*}
P\left(\bigcup_{i=1}^{k+1} A_i\right) &= P\left(\bigcup_{i=1}^k A_i \cup A_{k+1}\right) \\
&= P\left(\bigcup_{i=1}^k A_i\right) + P(A_{k+1}) - P\left(\bigcup_{i=1}^k A_i A_{k+1}\right) \\
& \leq \sum_{i=1}^k P(A_i) + P(A_{k+1}) = \sum_{i=1}^{k+1} P(A_i)
\end{align*}

From the principle of mathematical induction, for all $n$, the inequality holds.

This complete the whole proof.$\square$
\end{proof}

\section{Supplementary Problem 1.1}

The event "At least one of the events A,B,C happens" correspond to the set $A \cup B \cup C$.

$$
P(A \cup B \cup C) = P(A) + P(B) + P(C) - P(AB) - P(AC) - P(BC) + P(ABC)
$$
And from $P(AB) = 0$ we derives $P(ABC) = P(AB \cap C) = 0$.

Thus, we obtain:
$$
P(A \cup B \cup C) = \frac{3}{4} - 0 - 0 - \frac{1}{8} + 0 = \frac{5}{8}
$$

\section{Supplementary Problem 1.2}

From the given information, we know that 
$$
P(AB) = P(\bar{A} \cap B) = P(\overline{A \cup B}) = 1 - P(A \cup B)
$$

From the formula $P(A \cup B) = P(A) + P(B) - P(AB)$, we have:
$$
P(A \cup B) + P(AB) = P(A) + P(B) = p + P(B)
$$

Thus we conclude that 
$$
P(B) = 1 - p
$$

\section{Page 27, Problem 12}

\begin{align*}
N &= \binom{52}{5} \cdot \binom{47}{5} \cdot \binom{42}{5} \cdot \binom{37}{5} \cdot \binom{32}{5}
\end{align*}

\section{Page 29, Problem 29}
Let 
\begin{itemize}
  \item event $A$ be $\{\text{the player is dealt 3 spades and two hearts, and he discards two spades}\}$.
  \item event $B$ be $\{\text{the player gets two more spades}\}$
\end{itemize}
Then we have:
\begin{align*}
P(B|A) = \frac{\binom{10}{2}}{\binom{47}{2}} = \frac{45}{1081} 
\end{align*}

\section{Supplementary Problem 2.1}

The problem is as follows: 

Choose $2r$ shoes from $n$ pairs of shoes. Compute the probability of the following events:

\subsection{a) No Matching Pair}

\paragraph{Method 1}

$$
|\Omega| = \binom{2n}{2r}
$$

and 

\begin{align*}
|A| &= |\Omega| - |\{\text{there is at least 1 matching pair}\}| + \cdots + (-1)^r |\{\text{there are at least r matching pairs}\}| \\
    &= \sum_{i=0}^r (-1)^i \binom{n}{i} \binom{2(n-i)}{2r-2i}
\end{align*}

The probability of the event "No Matching Pair" is 
\begin{align*}
P(\text{No Matching Pair}) &= \frac{|A|}{|\Omega|} = \frac{\sum_{i=0}^r (-1)^i \binom{n}{i} \binom{2(n-i)}{2r-2i}}{\binom{2n}{2r}}
\end{align*}

\paragraph{Method 2} 

We choose $2r$ shoes from $n$ pairs of shoes, with each pair of them we choose the left shoe or the right shoe. Thus, we have:
\begin{align*}
|A| &= \binom{n}{2r} \cdot 2^{2r}
\end{align*}

Thus, the probability of the event "No Matching Pair" is
\begin{align*}
P(\text{No Matching Pair}) &= \frac{|A|}{|\Omega|} = \frac{\binom{n}{2r} \cdot 2^{2r}}{\binom{2n}{2r}}
\end{align*}

It can be proven that the two methods are equivalent.

\subsection{b) There is exactly one matching pair} 

We choose one pair of shoes from $n$ pairs of shoes, and then we choose $2r-2$ shoes from the remaining $n-1$ pairs of shoes, with each pair of them we choose the left shoe or the right shoe. Thus, we have:
$$
P(\text{Exactly one matching pair}) = \frac{\binom{n}{1} \cdot \binom{n-1}{2r-2} \cdot 2^{2r-2}}{\binom{2n}{2r}}
$$

\subsection{c) There are exactly r matching pairs}

Similar to the method in part b), we have:
$$
P(\text{Exactly r matching pairs}) = \frac{\binom{n}{r}}{\binom{2n}{2r}}
$$

\begin{remark} 
It can be proven that for $k$ matching pairs, the result is 
$$
P(\text{Exactly $k$ matching pairs}) = \frac{\binom{n}{k} \cdot \binom{n-k}{2(r-k)}\cdot 2^{2(r-k)}}{\binom{2n}{2r}}
$$
\end{remark}


\end{document}